\clearpage\begingroup\SourceLeaf{395b}\ArithmeticHead{64}{L}{STÄCKEL, GAUSS ALS GEOMETER.}
\renewcommand{\thefootnote}{\arabic{footnote})}\setcounter{footnote}{1}
\begin{center}
\textbf{18.}\par\smallskip
\textls[100]{Benutzung der komplexen Grössen für geometrische Untersuchungen.}
\end{center}
\noindent\hspace*{\parindent}Es ist eine merkwürdige Tatsache, dass \textsc{Gauss} fast überall, wo er mit
seinen Forschungen einsetzte, auf die komplexen Grössen stiess. Gilt das,
wie wir gesehen haben, für die Algebra, die Funktionentheorie und die Arithmetik,
so ist es nicht minder richtig für die Geometrie selbst.
\clearpage\SourceLeaf{396a}\ArithmeticHead{65}{R}{DIE KOMPLEXEN GRÖSSEN IN IHRER BEZIEHUNG ZUR GEOMETRIE.}\setcounter{footnote}{0}
\noindent\hspace*{\parindent}Die unmittelbare Anwendung der geometrischen Versinnlichung der komplexen
Grössen auf das Dreieck, das Viereck, den Kreis, die Kegelschnitte,
die Kugel ist ein Gegenstand, mit dem sich \textsc{Gauss} sein ganzes Leben lang
immer wieder beschäftigt hat; ja er hat diese Art der Behandlung geometrischer
Probleme als »eine ihm eigentümliche Methode« bezeichnet (Brief an
\textsc{Schumacher} vom 12. Mai 1843, W. VIII, S. 295). Die betreffenden Untersuchungen
werden im vierten Abschnitt dieses Aufsatzes im Zusammenhang
mit den Arbeiten zur elementaren und analytischen Geometrie ausführlich
dargestellt werden.
\par
Ausserdem ist die konforme Abbildung krummer Flächen zu erwähnen.
Allerdings hat \textsc{Gauss} in der Kopenhagener Preisschrift vom Jahre 1822 sich
bezüglich der geometrischen Versinnlichung auf Andeutungen beschränkt, die
kaum über das hinausgehen, was man in seiner Dissertation lesen kann. Im
übrigen sei auf die Darstellung im fünften Abschnitt dieses Aufsatzes verwiesen.
\clearpage\endgroup

